\chapter{2021年全国乙卷（文）}

\section{单选题}

\begin{problem}
已知全集 \(U=\{1,2,3,4,5\}\)，集合 \(M=\{1,2\}\)，\(N=\{3,4\}\)，则 \(C_U\left(M\cup N\right)=\)（\quad）

\choices
    {\(\{5\}\)}
    {\(\{1,2\}\)}
    {\(\{3,4\}\)}
    {\(\{1,2,3,4\}\)}
\end{problem}

\begin{problem}
设 \(\i z=4+3\i\)，则 \(z=\)（\quad）

\choices
    {\(-3-4\i\)}
    {\(-3+4\i\)}
    {\(3-4\i\)}
    {\(3+4\i\)}
\end{problem}

\begin{problem}
已知命题 \(p:\exists x\in\R\)，\(\sin x<1\)；命题 \(q:\forall x\in\R\)，\(\e^{|x|}\ge1\)，则下列命题中为真命题的是（\quad）

\choices
    {\(p\wedge q\)}
    {\(\neg p\wedge q\)}
    {\(p\wedge \neg q\)}
    {\(\neg\left(p\vee q\right)\)}
\end{problem}

\begin{problem}
函数 \(f(x)=\sin\dfrac{x}{3}+\cos\dfrac{x}{3}\) 的最小正周期和最大值分别是（\quad）

\choices
    {\(3\pi\) 和 \(\sqrt2\)}
    {\(3\pi\) 和 \(2\)}
    {\(6\pi\) 和 \(\sqrt2\)}
    {\(6\pi\) 和 \(2\)}
\end{problem}

\begin{problem}
若 \(x\)，\(y\) 满足约束条件
\(\begin{cases}
x+y\ge4,\\
x-y\le2,\\
y\le3,
\end{cases}\)
则 \(z=3x+y\) 的最小值为（\quad）

\choices
    {\(18\)}
    {\(10\)}
    {\(6\)}
    {\(4\)}
\end{problem}

\begin{problem}
\(\cos^2\dfrac{\pi}{12}-\cos^2\dfrac{5\pi}{12}=\)（\quad）

\choices
    {\(\dfrac{1}{2}\)}
    {\(\dfrac{\sqrt3}{3}\)}
    {\(\dfrac{\sqrt2}{2}\)}
    {\(\dfrac{\sqrt3}{2}\)}
\end{problem}

\begin{problem}
在区间 \(\left(0,\dfrac{1}{2}\right)\) 随机取 \(1\) 个数，则取到的数小于 \(\dfrac{1}{3}\) 的概率为（\quad）

\choices
    {\(\dfrac{3}{4}\)}
    {\(\dfrac{2}{3}\)}
    {\(\dfrac{1}{3}\)}
    {\(\dfrac{1}{6}\)}
\end{problem}

\begin{problem}
下列函数中最小值为 \(4\) 的是（\quad）

\choices
    {\(y=x^2+2x+4\)}
    {\(y=|\sin x|+\dfrac{4}{|\sin x|}\)}
    {\(y=2^x+2^{2-x}\)}
    {\(y=\ln x+\dfrac{4}{\ln x}\)}
\end{problem}

\begin{problem}
设函数 \(f(x)=\dfrac{1-x}{1+x}\)，则下列函数中为奇函数的是（\quad）

\choices
    {\(f(x-1)-1\)}
    {\(f(x-1)+1\)}
    {\(f(x+1)-1\)}
    {\(f(x+1)+1\)}
\end{problem}

\begin{problem}
在正方体 \(ABCD-A_1B_1C_1D_1\) 中，\(P\) 为 \(B_1D_1\) 的中点，则直线 \(PB\) 与 \(AD_1\) 所成的角为（\quad）

\choices
    {\(\dfrac{\pi}{2}\)}
    {\(\dfrac{\pi}{3}\)}
    {\(\dfrac{\pi}{4}\)}
    {\(\dfrac{\pi}{6}\)}
\end{problem}

\begin{problem}
设 \(B\) 是椭圆 \(C:\dfrac{x^2}{5}+y^2=1\) 的上顶点，点 \(P\) 在 \(C\) 上，则 \(\lvert PB\rvert\) 的最大值为（\quad）

\choices
    {\(\dfrac{5}{2}\)}
    {\(\sqrt6\)}
    {\(\sqrt5\)}
    {2}
\end{problem}

\begin{problem}
设 \(a\ne0\)，若 \(x=a\) 是函数 \(f(x)=a\left(x-a\right)^2\left(x-b\right)\) 的极大值点，则（\quad）

\choices
    {\(a<b\)}
    {\(a>b\)}
    {\(ab<a^2\)}
    {\(ab>a^2\)}
\end{problem}

\section{填空题}

\begin{problem}
已知向量 \(\bs{a}=(2,5)\)，\(\bs{b}=(\lambda,4)\)，若 \(\bs{a}\parallel\bs{b}\)，则 \(\lambda=\)\fillinblank{}.
\end{problem}

\begin{problem}
双曲线 \(\dfrac{x^2}{4}-\dfrac{y^2}{5}=1\) 的右焦点到直线 \(x+2y-8=0\) 的距离为\fillinblank{}.
\end{problem}

\begin{problem}
记 \(\triangle ABC\) 的内角 \(A\)，\(B\)，\(C\) 的对边分别为 \(a\)，\(b\)，\(c\)，面积为 \(\sqrt3\)，\(B=60^\circ\)，\(a^2+c^2=3ac\)，则 \(b=\)\fillinblank{}.
\end{problem}

\begin{problem}
以图\circled{1}为正视图，在图\circled{2}，\circled{3}，\circled{4}，\circled{5}中选两个分别作为侧视图和俯视图，组成某个三棱锥的三视图，则所选侧视图和俯视图的编号依次为\fillinblank{}（写出符合要求的一组答案即可）.

\bitmapfigure[width=17cm]{img/2021/national_paper_B_liberal/q16_fig1.png}
\end{problem}

\section{解答题}

\begin{problem}
某厂研制了一种生产高精产品的设备，为检验新设备生产产品的某项指标有无提高，用一台旧设备和一台新设备各生产了 \(10\) 件产品，得到各件产品该项指标数据如下：

\begin{center}
\begin{tabular}{|c|*{10}{c|}}
\hline
旧设备 & \(9.8\) & \(10.3\) & \(10.0\) & \(10.2\) & \(9.9\) & \(9.8\) & \(10.0\) & \(10.1\) & \(10.2\) & \(9.7\) \\
\hline
新设备 & \(10.1\) & \(10.4\) & \(10.1\) & \(10.0\) & \(10.1\) & \(10.3\) & \(10.6\) & \(10.5\) & \(10.4\) & \(10.5\) \\
\hline
\end{tabular}
\end{center}

旧设备和新设备生产产品的该项指标的样本平均数分别记为 \(\bar{x}\) 和 \(\bar{y}\)，样本方差分别记为 \(s_1^2\) 和 \(s_2^2\).

\begin{enumerate}
\item 求 \(\bar{x}\)，\(\bar{y}\)，\(s_1^2\)，\(s_2^2\)；
\item 判断新设备生产产品的该项指标的均值较旧设备是否有显著提高（如果 \(\bar{y}-\bar{x}\ge2\sqrt{\dfrac{s_1^2+s_2^2}{10}}\)，则认为新设备生产产品的该项指标的均值较旧设备有显著提高，否则不认为有显著提高）.
\end{enumerate}
\end{problem}

\begin{problem}
如图，四棱锥 \(P-ABCD\) 的底面是矩形，\(PD\perp\) 底面 \(ABCD\)，\(M\) 为 \(BC\) 的中点，且 \(PB\perp AM\).

\begin{enumerate}
\item 证明：平面 \(PAM\perp\) 平面 \(PBD\)；
\item 若 \(PD=DC=1\)，求四棱锥 \(P-ABCD\) 的体积.
\end{enumerate}

\bitmapfigure[width=0.42\linewidth]{img/2021/national_paper_B_liberal/q18_fig1.png}
\end{problem}

\begin{problem}
设 \(\{a_n\}\) 是首项为 \(1\) 的等比数列，数列 \(\{b_n\}\) 满足 \(b_n=\dfrac{na_n}{3}\). 已知 \(a_1\)，\(3a_2\)，\(9a_3\) 成等差数列.

\begin{enumerate}
\item 求 \(\{a_n\}\) 和 \(\{b_n\}\) 的通项公式；
\item 记 \(S_n\) 和 \(T_n\) 分别为 \(\{a_n\}\) 和 \(\{b_n\}\) 的前 \(n\) 项和，证明：\(T_n<\dfrac{S_n}{2}\).
\end{enumerate}
\end{problem}

\begin{problem}
已知抛物线 \(C:y^2=2px\)（\(p>0\)）的焦点 \(F\) 到准线的距离为 \(2\).

\begin{enumerate}
\item 求 \(C\) 的方程；
\item 已知 \(O\) 为坐标原点，点 \(P\) 在 \(C\) 上，点 \(Q\) 满足 \(\overrightarrow{PQ}=9\overrightarrow{QF}\)，求直线 \(OQ\) 斜率的最大值.
\end{enumerate}
\end{problem}

\begin{problem}
已知函数 \(f(x)=x^3-x^2+ax+1\).

\begin{enumerate}
\item 讨论 \(f(x)\) 的单调性；
\item 求曲线 \(y=f(x)\) 过坐标原点的切线与曲线 \(y=f(x)\) 的公共点的坐标.
\end{enumerate}
\end{problem}

\section{选考题：解答题}

\begin{problem}
在直角坐标系 \(xOy\) 中，\(\odot C\) 的圆心为 \(C(2,1)\)，半径为 \(1\).

\begin{enumerate}
\item 写出 \(\odot C\) 的一个参数方程；
\item 过点 \(F(4,1)\) 作 \(\odot C\) 的两条切线. 以坐标原点为极点，\(x\) 轴正半轴为极轴建立坐标系，求这两条切线的极坐标方程.
\end{enumerate}
\end{problem}

\begin{problem}
已知函数 \(f(x)=|x-a|+|x+3|\).

\begin{enumerate}
\item 当 \(a=1\) 时，求不等式 \(f(x)\ge6\) 的解集；
\item 若 \(f(x)>-a\)，求 \(a\) 的取值范围.
\end{enumerate}
\end{problem}

